\documentclass[12pt]{article}
\usepackage[margin=1in]{geometry}
\usepackage{amsmath}
\usepackage{enumitem}
\usepackage{array}
\usepackage{xcolor}
\usepackage{listings}

\title{EC441 Worked Problems: DNS, HTTP, and QUIC}
\author{}
\date{}

\lstset{
    basicstyle=\ttfamily\small,
    breaklines=true,
    frame=single
}

\begin{document}

\maketitle

\section*{Problem 1: DNS Namespace}

Break down the DNS name:

\[
\texttt{www.eng.bu.edu.}
\]

Identify the host, subdomain, second-level domain, top-level domain, and root.

\subsection*{Solution}

DNS names are read from right to left.

\[
\texttt{www.eng.bu.edu.}
\]

\[
\begin{array}{c|c}
\text{Component} & \text{Meaning} \\
\hline
\texttt{.} & \text{Root} \\
\texttt{edu} & \text{Top-level domain} \\
\texttt{bu} & \text{Second-level domain} \\
\texttt{eng} & \text{Subdomain} \\
\texttt{www} & \text{Host}
\end{array}
\]

\[
\boxed{\texttt{www.eng.bu.edu.} \text{ is a hierarchical DNS name.}}
\]

\section*{Problem 2: DNS Resolution Chain}

A laptop wants the IP address for:

\[
\texttt{www.eng.bu.edu}
\]

Assume the recursive resolver has nothing cached. List the steps needed to resolve the name.

\subsection*{Solution}

The laptop first sends the query to a recursive resolver.

Then the recursive resolver walks the DNS hierarchy:

\begin{enumerate}
    \item Ask a root server: ``Who handles \texttt{edu}?''
    \item Ask an \texttt{edu} TLD server: ``Who handles \texttt{bu.edu}?''
    \item Ask a \texttt{bu.edu} authoritative server: ``Who handles \texttt{eng.bu.edu}?''
    \item Ask a \texttt{eng.bu.edu} authoritative server: ``What is the A record for \texttt{www.eng.bu.edu}?''
\end{enumerate}

The answer then returns back through the resolver to the laptop.

\[
\boxed{\text{The laptop asks the recursive resolver. The resolver does the walking.}}
\]

\section*{Problem 3: DNS TTL Calculation}

A DNS record has a TTL of:

\[
1800 \text{ seconds}
\]

A recursive resolver receives the answer at:

\[
t = 240 \text{ seconds}
\]

At what time does the cached record expire?

\subsection*{Solution}

The expiration time is:

\[
t_{\text{expire}} = t_{\text{received}} + TTL
\]

Substitute values:

\[
t_{\text{expire}} = 240 + 1800
\]

\[
t_{\text{expire}} = 2040
\]

\[
\boxed{2040 \text{ seconds}}
\]

After this time, the resolver must query again instead of using the cached value.

\section*{Problem 4: DNS Record Types}

Fill in the purpose of each DNS record type.

\[
\begin{array}{c|c}
\text{Record Type} & \text{Purpose} \\
\hline
A & ? \\
AAAA & ? \\
CNAME & ? \\
MX & ? \\
NS & ? \\
TXT & ?
\end{array}
\]

\subsection*{Solution}

\[
\begin{array}{c|p{9cm}}
\text{Record Type} & \text{Purpose} \\
\hline
A & Maps a name to an IPv4 address. \\
AAAA & Maps a name to an IPv6 address. \\
CNAME & Creates an alias from one name to another name. \\
MX & Lists mail exchange servers for a domain. \\
NS & Lists name servers for a zone. \\
TXT & Stores text data such as SPF, DKIM, or domain verification records.
\end{array}
\]

\[
\boxed{\text{Different DNS record types answer different kinds of questions.}}
\]

\section*{Problem 5: DNS Transport}

For each DNS transport, identify the port and main purpose.

\[
\begin{array}{c|c|c}
\text{Transport} & \text{Port} & \text{Use} \\
\hline
UDP & ? & ? \\
TCP & ? & ? \\
TLS & ? & ? \\
HTTPS & ? & ?
\end{array}
\]

\subsection*{Solution}

\[
\begin{array}{c|c|p{8cm}}
\text{Transport} & \text{Port} & \text{Use} \\
\hline
UDP & 53 & Default DNS transport for small queries and replies. \\
TCP & 53 & Large responses, zone transfers, and DNSSEC-related traffic. \\
TLS & 853 & DNS over TLS, used for privacy. \\
HTTPS & 443 & DNS over HTTPS, used for privacy and blending with web traffic.
\end{array}
\]

\[
\boxed{\text{Traditional DNS is usually UDP port 53.}}
\]

\section*{Problem 6: HTTP Request}

Write a valid HTTP/1.1 GET request for:

\[
\texttt{http://example.com/index.html}
\]

\subsection*{Solution}

A valid HTTP/1.1 request is:

\begin{lstlisting}
GET /index.html HTTP/1.1
Host: example.com
User-Agent: curl/8.4.0
Accept: */*
\end{lstlisting}

The request line contains:

\[
\text{Method} + \text{Path} + \text{Version}
\]

\[
\boxed{\texttt{GET /index.html HTTP/1.1}}
\]

\section*{Problem 7: HTTP Response Interpretation}

Given the response:

\begin{lstlisting}
HTTP/1.1 200 OK
Content-Type: text/html; charset=UTF-8
Content-Length: 1256
Connection: keep-alive
\end{lstlisting}

Answer:

\begin{enumerate}
    \item Was the request successful?
    \item What type of content was returned?
    \item How large is the response body?
    \item Will the TCP connection stay open?
\end{enumerate}

\subsection*{Solution}

\begin{enumerate}
    \item Yes. The status code is:
    \[
    200
    \]
    which means success.

    \item The content type is:
    \[
    \texttt{text/html}
    \]
    so the response contains HTML.

    \item The response body is:
    \[
    1256 \text{ bytes}
    \]

    \item Yes. The header:
    \[
    \texttt{Connection: keep-alive}
    \]
    means the TCP connection may be reused.
\end{enumerate}

\[
\boxed{\text{This is a successful HTTP response carrying an HTML document.}}
\]

\section*{Problem 8: HTTP Status Code Classes}

Classify the following HTTP status codes.

\[
\begin{array}{c|c}
\text{Status Code} & \text{Class} \\
\hline
100 & ? \\
200 & ? \\
301 & ? \\
404 & ? \\
500 & ?
\end{array}
\]

\subsection*{Solution}

\[
\begin{array}{c|c}
\text{Status Code} & \text{Class} \\
\hline
100 & \text{Informational} \\
200 & \text{Success} \\
301 & \text{Redirect} \\
404 & \text{Client error} \\
500 & \text{Server error}
\end{array}
\]

\[
\boxed{
\begin{array}{l}
1xx = \text{Informational} \\
2xx = \text{Success} \\
3xx = \text{Redirect} \\
4xx = \text{Client error} \\
5xx = \text{Server error}
\end{array}
}
\]

\section*{Problem 9: HTTP Statelessness}

HTTP is stateless, but websites still keep users logged in. Explain how this is possible.

\subsection*{Solution}

HTTP is stateless because each request stands alone. The protocol itself does not remember previous requests.

Login state is usually handled using cookies.

The process is:

\begin{enumerate}
    \item User logs in.
    \item Server creates a session.
    \item Server sends a cookie to the browser.
    \item Browser stores the cookie.
    \item Browser sends the cookie on future requests.
    \item Server looks up the session associated with the cookie.
\end{enumerate}

\[
\boxed{\text{The protocol is stateless, but cookies carry session identity.}}
\]

\section*{Problem 10: HTTP/1.1 Head-of-Line Blocking}

A client sends three HTTP/1.1 pipelined requests:

\[
R_1,\ R_2,\ R_3
\]

The server can generate responses for \(R_2\) and \(R_3\) quickly, but \(R_1\) is slow.

What happens?

\subsection*{Solution}

In HTTP/1.1 pipelining, responses must return in the same order as requests.

So even if \(R_2\) and \(R_3\) are ready, they cannot be sent before \(R_1\).

\[
R_1 \text{ slow} \Rightarrow R_2 \text{ waits} \Rightarrow R_3 \text{ waits}
\]

\[
\boxed{\text{One slow response blocks all later responses.}}
\]

This is application-layer head-of-line blocking.

\section*{Problem 11: HTTP/2 Multiplexing}

Explain how HTTP/2 solves HTTP/1.1 application-layer head-of-line blocking.

\subsection*{Solution}

HTTP/2 uses binary framing and multiplexing.

Instead of sending whole responses strictly in order, HTTP/2 splits messages into frames.

Frames from different streams can be interleaved:

\[
S_1,\ S_2,\ S_3
\]

can share the same TCP connection.

For example:

\[
S_1 \text{ frame},\ S_2 \text{ frame},\ S_3 \text{ frame},\ S_1 \text{ frame}
\]

This means a slow response does not block unrelated responses at the HTTP layer.

\[
\boxed{\text{HTTP/2 solves application-layer HOL blocking using multiplexed streams.}}
\]

\section*{Problem 12: HTTP/2 Remaining Problem}

HTTP/2 multiplexes many streams over one TCP connection. What head-of-line blocking problem remains?

\subsection*{Solution}

HTTP/2 still runs over TCP.

TCP delivers bytes in order. If one TCP segment is lost, TCP cannot deliver later bytes to the application until the missing segment is retransmitted.

Therefore, all HTTP/2 streams sharing that TCP connection can stall.

\[
\text{Lost TCP segment} \Rightarrow \text{TCP waits} \Rightarrow \text{all HTTP/2 streams wait}
\]

\[
\boxed{\text{HTTP/2 still suffers from TCP-layer head-of-line blocking.}}
\]

\section*{Problem 13: HTTP/3 and QUIC}

Why does HTTP/3 run over QUIC instead of TCP?

\subsection*{Solution}

HTTP/3 uses QUIC, which is a reliable transport protocol built on top of UDP.

QUIC provides independent streams. If one packet is lost in one stream, other streams can continue.

\[
\text{Loss in Stream 1} \not\Rightarrow \text{Stream 2 blocked}
\]

QUIC also includes:

\begin{itemize}
    \item Built-in TLS 1.3
    \item Faster connection setup
    \item 0-RTT resumption
    \item Connection migration
\end{itemize}

\[
\boxed{\text{HTTP/3 uses QUIC to remove TCP-layer head-of-line blocking.}}
\]

\section*{Problem 14: HTTP Version Timeline}

Fill in the missing entries.

\[
\begin{array}{c|c|c|c}
\text{Version} & \text{Transport} & \text{Wire Format} & \text{Main Improvement} \\
\hline
\text{HTTP/1.0} & ? & ? & ? \\
\text{HTTP/1.1} & ? & ? & ? \\
\text{HTTP/2} & ? & ? & ? \\
\text{HTTP/3} & ? & ? & ?
\end{array}
\]

\subsection*{Solution}

\[
\begin{array}{c|c|c|p{5cm}}
\text{Version} & \text{Transport} & \text{Wire Format} & \text{Main Improvement} \\
\hline
\text{HTTP/1.0} & \text{TCP} & \text{text} & \text{Basic HTTP request/response} \\
\text{HTTP/1.1} & \text{TCP} & \text{text} & \text{Keep-alive and Host header} \\
\text{HTTP/2} & \text{TCP} & \text{binary} & \text{Multiplexing and HPACK} \\
\text{HTTP/3} & \text{QUIC over UDP} & \text{binary} & \text{Removes TCP HOL blocking and builds in TLS}
\end{array}
\]

\[
\boxed{\text{HTTP semantics stay mostly stable while transport and wire format evolve.}}
\]

\section*{Problem 15: Browser Loading a Web Page}

A browser loads one HTML page that references:

\[
4 \text{ CSS files},\quad 8 \text{ JavaScript files},\quad 20 \text{ images}
\]

How many total HTTP requests are needed, assuming no cache hits?

\subsection*{Solution}

The browser first fetches the HTML page:

\[
1 \text{ HTML request}
\]

Then it fetches each referenced resource:

\[
4 + 8 + 20 = 32
\]

Total requests:

\[
1 + 32 = 33
\]

\[
\boxed{33 \text{ HTTP requests}}
\]

\section*{Problem 16: MIME Types}

Match each MIME type to its content.

\[
\begin{array}{c|c}
\text{MIME Type} & \text{Content} \\
\hline
\texttt{text/html} & ? \\
\texttt{text/css} & ? \\
\texttt{application/json} & ? \\
\texttt{image/png} & ? \\
\texttt{video/mp4} & ?
\end{array}
\]

\subsection*{Solution}

\[
\begin{array}{c|c}
\text{MIME Type} & \text{Content} \\
\hline
\texttt{text/html} & \text{HTML document} \\
\texttt{text/css} & \text{Stylesheet} \\
\texttt{application/json} & \text{JSON data} \\
\texttt{image/png} & \text{PNG image} \\
\texttt{video/mp4} & \text{MP4 video}
\end{array}
\]

\[
\boxed{\text{The Content-Type header tells the browser how to interpret bytes.}}
\]

\section*{Problem 17: DNS vs HTTP}

Which happens first when visiting:

\[
\texttt{http://example.com}
\]

DNS or HTTP? Explain.

\subsection*{Solution}

DNS happens first.

The browser needs an IP address before it can open a transport connection to the web server.

The order is:

\begin{enumerate}
    \item Resolve \texttt{example.com} using DNS.
    \item Open a TCP connection to the returned IP address.
    \item Send an HTTP request.
    \item Receive an HTTP response.
\end{enumerate}

\[
\boxed{\text{DNS finds the server. HTTP requests the resource.}}
\]

\section*{Problem 18: WebSocket}

Why does WebSocket start as HTTP and then switch protocols?

\subsection*{Solution}

HTTP is request/response, which is not ideal for applications like:

\[
\text{chat, live dashboards, collaborative editing, multiplayer games}
\]

These applications need both sides to send data whenever events happen.

WebSocket begins with an HTTP upgrade request:

\begin{lstlisting}
GET /chat HTTP/1.1
Host: example.com
Upgrade: websocket
Connection: Upgrade
\end{lstlisting}

After the server responds with:

\[
101 \text{ Switching Protocols}
\]

the connection becomes a persistent, bidirectional channel.

\[
\boxed{\text{WebSocket upgrades HTTP into a full-duplex communication channel.}}
\]

\section*{Final Summary}

\[
\boxed{
\begin{array}{l}
\text{DNS maps human-readable names to IP addresses.} \\
\text{HTTP fetches resources from servers.} \\
\text{HTTP/1.1 is text-based and uses persistent TCP connections.} \\
\text{HTTP/2 adds binary framing and multiplexing.} \\
\text{HTTP/3 uses QUIC over UDP to remove TCP head-of-line blocking.} \\
\text{Browsers coordinate many HTTP requests to build one web page.}
\end{array}
}
\]

\end{document}